\documentclass[11pt]{article}

\usepackage[margin=1in]{geometry}
\usepackage{amsmath, amssymb, amsfonts}
\usepackage{booktabs}
\usepackage{enumitem}
\usepackage{graphicx}
\usepackage{hyperref}
\graphicspath{{../../artifacts/}}

\title{Classifier Ladder and Contracts}
\author{kinematic-classifier-sandbox}
\date{\today}

\begin{document}

\maketitle

\section{Purpose}

This document explains the implementation-linked classifier ladder in the
repository and the contract surfaces that let different method families share a
common evaluation path. The point of the ladder is not that every method should
win every benchmark. The point is that each level adds a distinct form of
evidence and makes a distinct claim about what the repo is proving. In the
canonical repo story, this document is the \textbf{Classifier / Filter Ladder}
pillar: it explains how evidence providers answer the common question ``what is
the next \(\ell_t(\cdot)\)?'' before the evaluation layer decides whether the
result deserves promotion.

\section{Role In The Repo Story}

The ladder is the algorithm spine between the selected corpus and the posterior
evaluation layer:
\begin{equation}
\text{Validated Corpus}
\rightarrow
\text{Evidence Provider}
\rightarrow
\ell_t(\cdot)
\rightarrow
\text{Posterior Update}
\rightarrow
\text{Evaluation / Promotion}.
\end{equation}
The ladder claim is stricter than ``more complexity is better.'' The claim is:
\begin{quote}
\textit{each rung adds one new evidence capability and must be justified by a
failure mode of the previous rung.}
\end{quote}
That is why the ladder is paired with the 1D witness suite. The witness
problems are not miscellaneous toys; they are the smallest controlled studies
that justify each rung before 3D lift.

\section{Problem And Notation}

Across the ladder, the common problem is to infer a class or mode label from a
measurement history \(z_{1:t}\). The most important shared symbols are:
\begin{itemize}[leftmargin=1.5em]
    \item \(z_t\): observation at time \(t\),
    \item \(c\): class label in non-switching methods,
    \item \(s\): mode label in switching methods,
    \item \(p_t(c)\): posterior over classes after time \(t\),
    \item \(L_t(c)\): class-conditioned likelihood term,
    \item \(\phi_t\): feature vector extracted from a trajectory prefix or window,
    \item \(\ell_t(c)\): unnormalized log score before normalization.
\end{itemize}

Every method in the ladder is therefore an answer to:
\begin{equation}
    \text{How should } \ell_t(\cdot) \text{ be constructed from } z_{1:t}\text{?}
\end{equation}
Once that log-score vector is available, the shared normalization surface is
\begin{equation}
    p_t(c) = \frac{\exp(\ell_t(c))}{\sum_j \exp(\ell_t(j))}.
\end{equation}

The current ladder uses the following evidence contract:
\begin{center}
\small
\begin{tabular}{p{0.18\textwidth}p{0.36\textwidth}p{0.34\textwidth}}
\toprule
Method & Evidence input \(y_t\) & Evidence term \(\ell_t(c)\) \\
\midrule
pointwise & current observation or local feature & \(\log p(y_t \mid c)\) \\
windowed & fixed-window feature vector \(\phi_t\) & \(\log p(\phi_t \mid c)\) \\
sequential Bayes & recursively emitted evidence stream & accumulated per-step log evidence \\
Kalman bank & innovation residual \(\nu_{t,c}\) & \(\log \mathcal{N}(\nu_{t,c}; 0, S_{t,c})\) \\
transition model & evidence plus transition prior & \(\ell_t(c) + \log \sum_j T_{jc}p_{t-1}(j)\) \\
\bottomrule
\end{tabular}
\end{center}

The pointwise rung is the pure local-evidence baseline. The windowed rung
aggregates a short history into a feature vector. The sequential Bayes rung
adds recursive memory. The Kalman bank replaces hand-built evidence with
model-based innovation likelihoods. The transition rung adds explicit switching
structure for mode changes.

\section{Common Assumptions}

The ladder shares a few numerical and methodological assumptions:
\begin{itemize}[leftmargin=1.5em]
    \item additive log evidence is preferred to direct probability products,
    \item posterior state from the previous step is the recursive memory surface,
    \item higher-complexity methods must be justified by a failure mode of
    lower-complexity methods,
    \item shared evaluation consumes standardized output rows rather than
    method-specific internal states.
\end{itemize}

\section{Ladder Structure}

The current ladder is:
\begin{enumerate}[leftmargin=1.8em]
    \item \texttt{pointwise\_baseline.py}
    \item \texttt{windowed\_baseline.py}
    \item \texttt{sequential\_bayes\_accumulator.py}
    \item \texttt{kalman\_filter\_bank.py}
    \item \texttt{transition\_matrix\_accumulator.py}
\end{enumerate}

Each level adds a specific capability:
\begin{itemize}[leftmargin=1.5em]
    \item \texttt{pointwise}: one-step evidence from direct observations
    \item \texttt{windowed}: short-history summary evidence
    \item \texttt{state\_estimate}: covariance-aware evidence from provided
    filtered states
    \item \texttt{accumulator}: recursive cumulative evidence
    \item \texttt{kalman}: state-space prediction plus innovation likelihoods
    \item \texttt{transition\_matrix}: explicit switching pressure across modes
\end{itemize}

The current reader-facing ladder can therefore be summarized as:
\begin{center}
\small
\begin{tabular}{p{0.12\textwidth}p{0.17\textwidth}p{0.20\textwidth}p{0.20\textwidth}p{0.18\textwidth}}
\toprule
Rung & Algorithm & Adds & Failure addressed & 1D witness \\
\midrule
0 & pointwise & local likelihood baseline & no audited local baseline & pointwise overlap \\
1 & windowed & short local history & outlier and local-noise fragility & windowed outlier/extrema \\
2 & sequential Bayes & recursive memory & pointwise history blindness & sequential history \\
3 & Kalman bank & dynamics-conditioned innovations & endpoint ambiguity under irregular timing & Kalman endpoint match \\
4 & transition matrix & explicit mode switching & static-class assumption & transition switching \\
5 & IMM proof & switching-aware state inference & demonstrated switching failures & advanced 1D switching witness \\
6 & particle filter bank & sampled nonlinear / non-Gaussian evidence & linear-Gaussian assumptions fail under drag, outliers, or mean reversion & nonlinear drag and OU witnesses \\
7 & RBPF & sampled latent mode path plus conditional Kalman state & mixed discrete/continuous latent structure & latent maneuver onset witness \\
\bottomrule
\end{tabular}
\end{center}
This table is the intended front-door interpretation of the deeper derivations
that follow.

\section{Pointwise Evidence Baseline}

\texttt{pointwise\_baseline.py} is the direct-evidence baseline. It answers the
question:
\begin{equation}
    \text{How far can pure observation evidence go without temporal memory?}
\end{equation}

For class \(c\) with mean \(\mu_c\), standard deviation \(\sigma_c\), and
recursive prior \(p_{t-1}(c)\), the implementation computes
\begin{align}
    \ell_t(c)
    &= \log p_{t-1}(c)
    - \frac{1}{2}\left[
        \log\!\left(2\pi \sigma_c^2\right)
        + \frac{(z_t - \mu_c)^2}{\sigma_c^2}
    \right], \\
    p_t(c)
    &= \frac{\exp(\ell_t(c))}{\sum_j \exp(\ell_t(j))}.
\end{align}
This is the exact update used by
\texttt{GaussianPointwiseClassifier.update}. The model carries no explicit
latent state; its entire memory is the posterior vector \(p_{t-1}\).
This rung is the right baseline when a single sample already carries enough
class signal and the corpus does not require temporal context.

\section{Windowed Evidence Baseline}

\texttt{windowed\_baseline.py} adds short-history feature rows and summary
statistics. This is the bridge from single-sample scoring to structured feature
evidence.

Related implementation surfaces:
\begin{itemize}[leftmargin=1.5em]
    \item \texttt{irregular\_window\_comparison.py}
    \item \texttt{common\_experiment\_classifier\_registry.py}
\end{itemize}

For each time \(t\), the code constructs a feature row
\begin{align}
    \phi_t
    &=
    \big[
        \texttt{running\_min},
        \texttt{running\_max},
        \texttt{running\_range},
        \notag \\
    &\qquad
        \texttt{robust\_min},
        \texttt{robust\_max},
        \texttt{trimmed\_range},
        \texttt{slope},
        \dots
    \big]_t
\end{align}
from the trajectory prefix or local window. The selected evidence family is
\begin{itemize}[leftmargin=1.5em]
    \item raw: \(\{\texttt{running\_min}, \texttt{running\_max}, \texttt{running\_range}, \texttt{slope}\}\),
    \item robust: \(\{\texttt{robust\_min}, \texttt{robust\_max}, \texttt{trimmed\_range}, \texttt{slope}\}\).
\end{itemize}
For each class \(c\),
\begin{equation}
    \ell_t(c)
    = \log p_{t-1}(c)
    + \sum_{f \in \mathcal{F}}
    \log \mathcal{N}\!\left(\phi_{t,f}; \mu_{c,f}, \sigma_{c,f}^2\right),
\end{equation}
where \(\mathcal{F}\) is the selected feature subset. The current code assumes
conditional independence between feature coordinates inside a row.
That conditional-independence approximation is the key modeling assumption of
the windowed rung.
Windowed evidence is the right compromise when short local context matters but
the method still needs a simple probabilistic interface.

\section{Covariance-Aware State Estimate Evidence}

Sometimes the input is not a raw observation stream but an already filtered
state estimate with covariance, for example
\begin{equation}
    (x_t, P_t) \quad \text{or} \quad (x_{t|t}, P_{t|t}).
\end{equation}
In that case the Kalman update has already been performed elsewhere, so the
right evidence atom is not a new innovation residual. The right rung is a
covariance-aware state-likelihood score.

The class-conditioned evidence can be written as
\begin{equation}
    \ell_t(c)
    =
    -\frac{1}{2}
    \left[
        (x_t - \mu_c)^\top \Sigma_{t,c}^{-1}(x_t - \mu_c)
        + \log |\Sigma_{t,c}|
        + d \log(2\pi)
    \right],
\end{equation}
where \(\Sigma_{t,c}\) is the supplied covariance or a class-conditioned
covariance model built from it. If the provided covariance is class-neutral,
then \(\Sigma_{t,c}\) can be taken as the supplied covariance plus any
class-specific process or measurement inflation used by the benchmark.

This rung is the right choice when:
\begin{itemize}[leftmargin=1.5em]
    \item the tracker already gives a state estimate and covariance,
    \item the task is to classify trajectories from filtered states rather than
    to re-run tracking,
    \item the evidence should measure class fit to the filtered state, not the
    raw measurement innovation.
\end{itemize}
In other words, the ladder step is \emph{covariance-aware state evidence}, not
full Kalman filtering.

\section{Sequential Bayes Accumulator}

\texttt{sequential\_bayes\_accumulator.py} turns class evidence into recursive
posterior updates and is the cleanest additive log-likelihood ladder step before
latent-state filtering.

When a likelihood stream \(L_t(c)\) is supplied directly, the accumulator uses
\begin{align}
    \ell_t(c)
    &= \lambda \log p_{t-1}(c) + \log L_t(c), \\
    p_t(c)
    &= \frac{\exp(\ell_t(c))}{\sum_j \exp(\ell_t(j))},
\end{align}
where \(\lambda > 0\) is the forgetting factor in
\texttt{SequentialBayesAccumulator.update\_with\_likelihoods}. When the
accumulator is used in its direct Gaussian form, the likelihood itself is
\begin{equation}
    L_t(c) = \mathcal{N}(z_t; \mu_c, \sigma_c^2).
\end{equation}

The implementation also exposes a confidence gate. Let
\(c^\star_t = \arg\max_c p_t(c)\). The reported label is
\begin{equation}
    \hat{y}_t =
    \begin{cases}
        c^\star_t, & p_t(c^\star_t) \ge \tau, \\
        \texttt{unknown}, & p_t(c^\star_t) < \tau,
    \end{cases}
\end{equation}
with threshold \(\tau\) supplied as \texttt{confidence\_threshold}. That
abstention rule is why the accumulator artifacts can separate evidence weakness
from plain misclassification.
This rung is the simplest place where posterior quality becomes a first-class
question: a high posterior is only useful if it is also well calibrated.

\section{Kalman Innovation Bank}

\texttt{kalman\_filter\_bank.py} is the first full model-based recursive
classifier in the ladder. Its key ingredients are:
\begin{itemize}[leftmargin=1.5em]
    \item state prediction
    \item measurement innovation
    \item innovation likelihood
    \item model comparison across motion hypotheses
\end{itemize}

Related modules:
\begin{itemize}[leftmargin=1.5em]
    \item \texttt{kalman\_variant\_comparison.py}
    \item \texttt{kalman\_observable\_comparison.py}
    \item \texttt{velocity\_aided\_kalman\_comparison.py}
\end{itemize}

For each model \(m\), the code performs the standard linear-Gaussian predict and
update recursion under a class-conditioned motion model \cite{kalman1960}.
update steps
\begin{align}
    \hat{x}_{t|t-1}^{(m)} &= F_t^{(m)} x_{t-1|t-1}^{(m)}, \\
    P_{t|t-1}^{(m)} &= F_t^{(m)} P_{t-1|t-1}^{(m)} {F_t^{(m)}}^\top + Q_t^{(m)}, \\
    \nu_t^{(m)} &= z_t - H_t^{(m)} \hat{x}_{t|t-1}^{(m)}, \\
    S_t^{(m)} &= H_t^{(m)} P_{t|t-1}^{(m)} {H_t^{(m)}}^\top + R_t^{(m)}.
\end{align}
The classification evidence is the innovation likelihood
\begin{equation}
    \log L_t(m)
    = -\frac{1}{2}\left[
        \log\!\left(2\pi S_t^{(m)}\right)
        + \frac{(\nu_t^{(m)})^2}{S_t^{(m)}}
    \right],
\end{equation}
which is exactly what \texttt{\_innovation\_log\_likelihood} computes. The
posterior over models is then updated by
\begin{equation}
    \log \tilde{p}_t(m) = \log p_{t-1}(m) + \log L_t(m),
\end{equation}
followed by log-sum-exp normalization.

The implementation also includes two robustness mechanisms that materially
change benchmark behavior:
\begin{itemize}[leftmargin=1.5em]
    \item adaptive measurement inflation when normalized innovation squared is
    above a gate,
    \item smoothed process-noise scaling when the innovation remains unusually
    large.
\end{itemize}
Those policies are what the variant-comparison artifacts are actually testing.
This rung is the most model-driven one in the ladder. It is strongest when the
class hypothesis really implies a state-space trajectory and weakest when the
corpus is better described by local features alone.

\section{Transition-Aware Accumulation}

\texttt{transition\_matrix\_accumulator.py} adds explicit switching structure to
the recursive evidence story and acts as the bridge toward stronger switching
models.

The transition-aware accumulator first extracts finite-difference kinematics
\((\hat{v}_t, \hat{a}_t)\) from the measurement history and scores each mode by
Gaussian emission terms over speed, signed acceleration, and absolute
acceleration:
\begin{equation}
    \log E_t(s)
    =
    \log \mathcal{N}(\hat{v}_t; \mu^{(v)}_s, {\sigma^{(v)}_s}^2)
    +
    \log \mathcal{N}(\hat{a}_t; \mu^{(a)}_s, {\sigma^{(a)}_s}^2)
    +
    \log \mathcal{N}(|\hat{a}_t|; \mu^{(|a|)}_s, {\sigma^{(|a|)}_s}^2).
\end{equation}

Without transition structure, the prior remains \(p_{t-1}(s)\). With a
transition matrix \(T\), the code propagates
\begin{equation}
    \bar{p}_t(s') = \sum_s p_{t-1}(s)\,T_{s,s'}.
\end{equation}
The mode posterior then becomes
\begin{equation}
    \log \tilde{p}_t(s) = \log \bar{p}_t(s) + \log E_t(s).
\end{equation}
This is the current pre-IMM rung: the repo is already testing whether explicit
mode persistence and switching probabilities explain the observed failure mode
before paying the complexity cost of a full interacting multiple-model design.

There is now a worked numeric example for this update path in the artifact
below.
\begin{center}
\footnotesize
\texttt{artifacts/transition\_matrix\_accumulator\_v1/transition\_matrix\_numeric\_walkthrough.md}
\end{center}
That walkthrough uses a real switching trajectory and shows:
\begin{itemize}[leftmargin=1.5em]
    \item the propagated prior at the first switched step,
    \item the per-mode emission terms,
    \item the resulting log numerators,
    \item and the posterior normalization over modes for several steps around
    the switch.
\end{itemize}

\begin{figure}[htbp]
    \centering
    \includegraphics[width=0.9\textwidth]{transition_matrix_accumulator_v1/transition_matrix_diagnostics.png}
    \caption{Transition-aware accumulation diagnostics.}
\end{figure}

\section{Rung Sufficiency And Escalation Equations}

The ladder is only useful if the repo can answer not just ``what is the next
rung?'' but ``has the current rung earned the right to stay?'' The
\texttt{rung\_sufficiency} package makes that decision explicit.

For study \(s\) at current rung \(r\), define the oracle gap
\begin{equation}
    g_{\text{oracle}}(s,r)
    =
    \max\big(0, A_{\text{oracle}}(s,r) - A_{\text{current}}(s,r)\big),
\end{equation}
the prior-fragility term
\begin{equation}
    f_{\text{prior}}(s,r)
    =
    1 - \text{prior\_sensitivity\_score}(s,r),
\end{equation}
and the measured next-rung gain
\begin{equation}
    \Delta_{r \rightarrow r^{+}}(s)
    =
    A_{\text{current}}(s,r^{+}) - A_{\text{current}}(s,r),
\end{equation}
where \(r^{+}\) is the configured next rung from the capability matrix.

The current implementation treats a study as learnable when
\begin{equation}
    \mathbf{1}_{\text{learnable}}(s,r)
    =
    \mathbf{1}\!\left[
        A_{\text{oracle}} \ge \tau_{\text{oracle}}
        \land
        \mathrm{overlap} \le \tau_{\text{overlap}}
        \land
        \mathrm{AUC}_{\text{pair}} \ge \tau_{\text{AUC}}
        \land
        m_{\text{post}} \ge \tau_{\text{margin}}
    \right],
\end{equation}
with the current default thresholds
\begin{equation}
    \tau_{\text{oracle}} = 0.85,\qquad
    \tau_{\text{overlap}} = 0.90,\qquad
    \tau_{\text{AUC}} = 0.70,\qquad
    \tau_{\text{margin}} = 0.05.
\end{equation}
This is the mathematical version of the repo rule that feature and corpus
limits should be identified before algorithm blame is assigned.

The promotion decision is then a gated piecewise rule. Let
\(\tau_{\text{prior}} = 0.12\) and \(\tau_{\Delta} = 0.05\). Then
\begin{equation}
    d(s,r)
    =
    \begin{cases}
        \texttt{revise\_corpus}, & \neg \text{can\_evaluate\_classifier}(s), \\
        \texttt{revise\_features}, & \text{feature gate fails}, \\
        \texttt{revise\_prior}, & f_{\text{prior}}(s,r) > \tau_{\text{prior}}, \\
        \texttt{feature\_limited}, & \text{learnability\_status}(s,r)=\texttt{feature\_limited}, \\
        \texttt{stay}, & r^{+} = \varnothing, \\
        \texttt{reject\_escalation}, & \neg \text{capability\_match}(r,r^{+}), \\
        \texttt{defer\_advanced}, & \Delta_{r \rightarrow r^{+}}(s)\ \text{is unavailable}, \\
        \texttt{promote}, & \Delta_{r \rightarrow r^{+}}(s) \ge \tau_{\Delta}, \\
        \texttt{stay}, & g_{\text{oracle}}(s,r) < \tau_{\text{gap}}, \\
        \texttt{defer\_advanced}, & \text{otherwise},
    \end{cases}
\end{equation}
where \(\tau_{\text{gap}} = 0.08\) in the default configuration. This is the
explicit ladder-sufficiency rule that stops the repo from escalating methods
only because they are available.

The switching witnesses are the cleanest examples because their measured gains
are hard-coded from dedicated benchmark outputs:
\begin{align}
    \Delta_{\text{kalman} \rightarrow \text{transition}}
    &=
    A_{\text{post-switch}}^{\text{transition}}
    -
    A_{\text{post-switch}}^{\text{kalman}}, \\
    \Delta_{\text{transition} \rightarrow \text{IMM}}
    &=
    A_{\text{post-switch}}^{\text{IMM}}
    -
    A_{\text{post-switch}}^{\text{transition}}.
\end{align}
That is why the current rung-sufficiency artifact can issue a real promotion
decision for transition-aware accumulation and a measured
\texttt{promote}/\texttt{defer\_advanced} decision for IMM.

\section{Calibration And Posterior Quality}

Posterior quality is not the same as accuracy. A method can be accurate but
poorly calibrated, or calibrated but not especially accurate. The repo tracks
both notions because a posterior is only useful if its confidence is
trustworthy.

The main calibration metrics are:
\begin{align}
    \mathrm{NLL}
    &=
    -\frac{1}{N}\sum_{i=1}^{N}\log p_i(y_i), \\
    \mathrm{Brier}
    &=
    \frac{1}{N}\sum_{i=1}^{N}\sum_{c\in\mathcal{C}}
    \left(p_i(c)-\mathbf{1}\{y_i=c\}\right)^2, \\
    H_i
    &=
    -\sum_{c\in\mathcal{C}} p_i(c)\log p_i(c), \\
    \mathrm{ECE}
    &=
    \sum_{b=1}^{B}\frac{|I_b|}{N}
    \left|\mathrm{acc}(I_b)-\mathrm{conf}(I_b)\right|.
\end{align}
Here \(H_i\) is posterior entropy, \(\mathrm{NLL}\) is negative
log-likelihood, \(\mathrm{Brier}\) is the squared-probability error, and
\(\mathrm{ECE}\) is expected calibration error computed over confidence bins.

Posterior quality is usually interpreted with three additional indicators:
\begin{itemize}[leftmargin=1.5em]
    \item high mean confidence with high error rate indicates overconfidence,
    \item low entropy with low error rate indicates sharp, well-supported
    posteriors,
    \item high disagreement across method families indicates that posterior
    quality depends strongly on the evidence model, not only on the class pair.
\end{itemize}

The current artifact bundle does not yet expose NLL, Brier, and ECE as primary
columns everywhere, so the repo currently uses confidence, entropy, margin,
confident-error rate, prior flips, and disagreement as the practical posterior
quality surface.

\subsection*{Current Repo Interpretation}

On the current 90-trajectory generated corpus slice, the posterior-quality
summary is:
\begin{itemize}[leftmargin=1.5em]
    \item \texttt{windowed\_raw}: accuracy \(0.978\), mean confidence \(0.996\),
    mean entropy \(0.010\), confident-error rate \(0.022\)
    \item \texttt{windowed\_robust}: accuracy \(0.956\), mean confidence \(1.000\),
    mean entropy \(0.000\), confident-error rate \(0.044\)
    \item \texttt{pointwise}: accuracy \(0.867\), mean confidence \(0.984\),
    mean entropy \(0.034\), confident-error rate \(0.133\)
    \item \texttt{sequential\_bayes}: accuracy \(0.822\), mean confidence \(0.967\),
    mean entropy \(0.053\), confident-error rate \(0.111\)
    \item \texttt{kalman\_bank}: accuracy \(0.678\), mean confidence \(0.924\),
    mean entropy \(0.198\), confident-error rate \(0.244\)
\end{itemize}
The most important reading is that the windowed methods are currently the best
posterior-quality performers on this corpus slice, while the Kalman bank is the
least accurate and most uncertain. The pointwise and sequential Bayes rungs
sit between those extremes.

The current disagreement rate is \(46/90\) trajectories, which means the method
families are not redundant. They are making meaningfully different claims about
the same trajectories, which is exactly what a ladder should expose.

\section{Implementation Contract}

The ladder only becomes a methodology framework when different methods can be
evaluated through a common output surface. The practical repo contract is:
\begin{equation}
    \text{method-specific state and evidence}
    \rightarrow
    \text{shared prediction rows}
    \rightarrow
    \text{shared metrics}.
\end{equation}
This contract is implemented across
\begin{itemize}[leftmargin=1.5em]
    \item \texttt{common\_experiment\_harness.py}
    \item \texttt{common\_experiment\_classifier\_registry.py}
    \item \texttt{shared\_evaluation.py}
    \item \texttt{contracts.py}
\end{itemize}
and the generic contract proof modules.

\section{Shared Ladder Benchmark And Dataset Comparison}

The comparison layer is implemented through:
\begin{itemize}[leftmargin=1.5em]
    \item \texttt{common\_dataset\_comparison.py}
    \item \texttt{technique\_comparison.py}
    \item \texttt{common\_1d\_study\_adapter.py}
\end{itemize}

These modules normalize different methods onto one benchmark surface. The
practical contract is:
\begin{equation}
    \text{method-specific evidence/history}
    \rightarrow
    \text{shared prediction rows}
    \rightarrow
    \text{shared metrics and artifact writers}.
\end{equation}
The comparison layer therefore does not assume that every method has the same
internal state. It only assumes that each method can emit compatible posterior,
prediction, and metric rows.

\section{Common Experiment Harness And Adapter Layer}

The shared study layer is implemented through:
\begin{itemize}[leftmargin=1.5em]
    \item \texttt{common\_experiment\_harness.py}
    \item \texttt{common\_experiment\_classifier\_registry.py}
    \item \texttt{shared\_evaluation.py}
\end{itemize}

These modules are the practical answer to:
\begin{equation}
    \text{different classifier family}
    \rightarrow
    \text{shared metrics and artifact layout}.
\end{equation}

At the schema level, the harness expects a method family to provide at least:
\begin{itemize}[leftmargin=1.5em]
    \item trajectory and scenario identifiers,
    \item true and predicted class labels,
    \item posterior probabilities,
    \item confidence values,
    \item optional evidence or filter diagnostics.
\end{itemize}
Everything downstream --- confusion matrices, prior-sensitivity views, feature
set studies, and class-pair studies --- is built on that shared surface.

\section{Common Output Schema And Contracts}

The explicit contract logic lives in:
\begin{itemize}[leftmargin=1.5em]
    \item \texttt{contracts.py}
    \item \texttt{generic\_inference\_contract.py}
    \item \texttt{generic\_classification\_evidence\_proof.py}
    \item \texttt{validation\_ladder.py}
\end{itemize}

These modules are the main proof that the repo is becoming a generic study
framework rather than a collection of disconnected method-specific writers.

The generic inference proof can be summarized as
\begin{equation}
    \{\log L_t(c)\}_{c \in \mathcal{C}}
    \xrightarrow{\text{posterior updater}}
    \{p_t(c)\}_{c \in \mathcal{C}},
\end{equation}
independent of whether \(\log L_t(c)\) came from direct measurement scoring,
windowed features, a recursive accumulator, or an innovation likelihood.

\section{Backend Adapter Proof}

The backend compatibility surfaces live in:
\begin{itemize}[leftmargin=1.5em]
    \item \texttt{generic\_filtering\_contract.py}
    \item \texttt{trajectory\_backend\_contract.py}
    \item \texttt{backend\_adapter\_proof.py}
\end{itemize}

The backend-adapter claim is deliberately narrower than ``all filters are the
same.'' The actual claim is:
\begin{align}
    \text{backend}
    \rightarrow
    \big(
        \text{state summary},
        \text{evidence summary},
        \text{diagnostics}
    \big)
    \rightarrow
    \text{shared evaluation contract}.
\end{align}
That is the contract the current IMM proof satisfies, and it now extends to the
shared PF branch, the RBPF branch, and the OU-style mean-reversion witness that
exercises a concrete stochastic pull-back-to-class dynamic inside the PF family.

\section{Artifact Map}

Primary artifacts for this layer:
\begin{itemize}[leftmargin=1.5em]
    \item \texttt{artifacts/common\_dataset\_comparison\_v1/}
    \item \texttt{artifacts/technique\_comparison\_v1/}
    \item \texttt{artifacts/generic\_inference\_contract/}
    \item \texttt{artifacts/filtering\_contract/}
    \item \texttt{artifacts/transition\_matrix\_accumulator\_v1/}
    \item \texttt{artifacts/advanced\_state\_inference\_contract/}
    \item \texttt{artifacts/advanced\_state\_inference\_v1/}
\end{itemize}

\section{References}
\begin{thebibliography}{9}
\bibitem{brier1950}
G. W. Brier.
\newblock Verification of forecasts expressed in terms of probability.
\newblock \emph{Monthly Weather Review}, 78(1):1--3, 1950.

\bibitem{niculescu2005}
A. Niculescu-Mizil and R. Caruana.
\newblock Predicting good probabilities with supervised learning.
\newblock In \emph{ICML}, 2005.

\bibitem{guo2017}
C. Guo, G. Pleiss, Y. Sun, and K. Q. Weinberger.
\newblock On calibration of modern neural networks.
\newblock In \emph{ICML}, 2017.

\bibitem{kalman1960}
R. E. Kalman.
\newblock A new approach to linear filtering and prediction problems.
\newblock \emph{Transactions of the ASME}, 1960.
\end{thebibliography}

\end{document}
